\section{Conclusão}

O desenvolvimento do \textit{OffScan} demonstrou a viabilidade de criar uma ferramenta unificada para testes ofensivos de redes utilizando Rust como linguagem principal. A ferramenta implementou com sucesso funcionalidades essenciais como flooding de pacotes, mapeamento de redes, escaneamento de portas e ataques a redes Wi-Fi, alcançando desempenho satisfatório em ambiente controlado.

Os resultados confirmaram as vantagens da escolha de Rust, particularmente em termos de segurança de memória e desempenho, enquanto a integração com C permitiu acesso a APIs de baixo nível do kernel. Embora existam oportunidades de otimização e expansão, o \textit{OffScan} representa uma contribuição válida para o ecossistema de ferramentas de segurança ofensiva, oferecendo uma base sólida para desenvolvimentos futuros.

Trabalhos subsequentes poderão focar na redução do \textit{overhead} de chamadas ao sistema e na implementação de novos tipos de ataques, transformando a ferramenta em uma solução mais completa para profissionais de segurança.